%%%%%%%%%%%%%%%%%%%%%%%%%%%%%%%%%%%%%%%%%
% Lachaise Assignment
% LaTeX Template
% Version 1.0 (26/6/2018)
%
% This template originates from:
% http://www.LaTeXTemplates.com
%
% Authors:
% Marion Lachaise & François Févotte
% Vel (vel@LaTeXTemplates.com)
%
% License:
% CC BY-NC-SA 3.0 (http://creativecommons.org/licenses/by-nc-sa/3.0/)
%
%%%%%%%%%%%%%%%%%%%%%%%%%%%%%%%%%%%%%%%%%

%----------------------------------------------------------------------------------------
%	PACKAGES AND OTHER DOCUMENT CONFIGURATIONS
%----------------------------------------------------------------------------------------

\documentclass{article}
\usepackage{amsthm}
\usepackage{listings}
\usepackage{xcolor}
\usepackage{graphicx}
\usepackage{cite}
\usepackage{longtable}
\usepackage{array}
\usepackage{booktabs}

\lstset{
	language=Python,
	basicstyle=\ttfamily\small,
	keywordstyle=\color{blue},
	stringstyle=\color{red},
	commentstyle=\color{purple},
	showstringspaces=false,
	numbers=left,
	numberstyle=\tiny\color{gray},
	breaklines=true,
	frame=single,
	captionpos=b
}
\newtheorem{definition}{Definition}[section]

%%%%%%%%%%%%%%%%%%%%%%%%%%%%%%%%%%%%%%%%%
% Lachaise Assignment
% Structure Specification File
% Version 1.0 (26/6/2018)
%
% This template originates from:
% http://www.LaTeXTemplates.com
%
% Authors:
% Marion Lachaise & François Févotte
% Vel (vel@LaTeXTemplates.com)
%
% License:
% CC BY-NC-SA 3.0 (http://creativecommons.org/licenses/by-nc-sa/3.0/)
% 
%%%%%%%%%%%%%%%%%%%%%%%%%%%%%%%%%%%%%%%%%

%----------------------------------------------------------------------------------------
%	PACKAGES AND OTHER DOCUMENT CONFIGURATIONS
%----------------------------------------------------------------------------------------

\usepackage{amsmath,amsfonts,stmaryrd,amssymb} % Math packages

\usepackage{enumerate} % Custom item numbers for enumerations

\usepackage[ruled]{algorithm2e} % Algorithms

\usepackage[framemethod=tikz]{mdframed} % Allows defining custom boxed/framed environments

\usepackage{listings} % File listings, with syntax highlighting
\lstset{
	basicstyle=\ttfamily, % Typeset listings in monospace font
}

%----------------------------------------------------------------------------------------
%	DOCUMENT MARGINS
%----------------------------------------------------------------------------------------

\usepackage{geometry} % Required for adjusting page dimensions and margins

\geometry{
	paper=a4paper, % Paper size, change to letterpaper for US letter size
	top=2.5cm, % Top margin
	bottom=3cm, % Bottom margin
	left=2.5cm, % Left margin
	right=2.5cm, % Right margin
	headheight=14pt, % Header height
	footskip=1.5cm, % Space from the bottom margin to the baseline of the footer
	headsep=1.2cm, % Space from the top margin to the baseline of the header
	%showframe, % Uncomment to show how the type block is set on the page
}

%----------------------------------------------------------------------------------------
%	FONTS
%----------------------------------------------------------------------------------------

\usepackage[utf8]{inputenc} % Required for inputting international characters
\usepackage[T1]{fontenc} % Output font encoding for international characters

\usepackage{XCharter} % Use the XCharter fonts

%----------------------------------------------------------------------------------------
%	COMMAND LINE ENVIRONMENT
%----------------------------------------------------------------------------------------

% Usage:
% \begin{commandline}
%	\begin{verbatim}
%		$ ls
%		
%		Applications	Desktop	...
%	\end{verbatim}
% \end{commandline}

\mdfdefinestyle{commandline}{
	leftmargin=10pt,
	rightmargin=10pt,
	innerleftmargin=15pt,
	middlelinecolor=black!50!white,
	middlelinewidth=2pt,
	frametitlerule=false,
	backgroundcolor=black!5!white,
	frametitle={Command Line},
	frametitlefont={\normalfont\sffamily\color{white}\hspace{-1em}},
	frametitlebackgroundcolor=black!50!white,
	nobreak,
}

% Define a custom environment for command-line snapshots
\newenvironment{commandline}{
	\medskip
	\begin{mdframed}[style=commandline]
}{
	\end{mdframed}
	\medskip
}

%----------------------------------------------------------------------------------------
%	FILE CONTENTS ENVIRONMENT
%----------------------------------------------------------------------------------------

% Usage:
% \begin{file}[optional filename, defaults to "File"]
%	File contents, for example, with a listings environment
% \end{file}

\mdfdefinestyle{file}{
	innertopmargin=1.6\baselineskip,
	innerbottommargin=0.8\baselineskip,
	topline=false, bottomline=false,
	leftline=false, rightline=false,
	leftmargin=2cm,
	rightmargin=2cm,
	singleextra={%
		\draw[fill=black!10!white](P)++(0,-1.2em)rectangle(P-|O);
		\node[anchor=north west]
		at(P-|O){\ttfamily\mdfilename};
		%
		\def\l{3em}
		\draw(O-|P)++(-\l,0)--++(\l,\l)--(P)--(P-|O)--(O)--cycle;
		\draw(O-|P)++(-\l,0)--++(0,\l)--++(\l,0);
	},
	nobreak,
}

% Define a custom environment for file contents
\newenvironment{file}[1][File]{ % Set the default filename to "File"
	\medskip
	\newcommand{\mdfilename}{#1}
	\begin{mdframed}[style=file]
}{
	\end{mdframed}
	\medskip
}

%----------------------------------------------------------------------------------------
%	NUMBERED QUESTIONS ENVIRONMENT
%----------------------------------------------------------------------------------------

% Usage:
% \begin{question}[optional title]
%	Question contents
% \end{question}

\mdfdefinestyle{question}{
	innertopmargin=1.2\baselineskip,
	innerbottommargin=0.8\baselineskip,
	roundcorner=5pt,
	nobreak,
	singleextra={%
		\draw(P-|O)node[xshift=1em,anchor=west,fill=white,draw,rounded corners=5pt]{%
		Question \theQuestion\questionTitle};
	},
}

\newcounter{Question} % Stores the current question number that gets iterated with each new question

% Define a custom environment for numbered questions
\newenvironment{question}[1][\unskip]{
	\bigskip
	\stepcounter{Question}
	\newcommand{\questionTitle}{~#1}
	\begin{mdframed}[style=question]
}{
	\end{mdframed}
	\medskip
}

%----------------------------------------------------------------------------------------
%	WARNING TEXT ENVIRONMENT
%----------------------------------------------------------------------------------------

% Usage:
% \begin{warn}[optional title, defaults to "Warning:"]
%	Contents
% \end{warn}

\mdfdefinestyle{warning}{
	topline=false, bottomline=false,
	leftline=false, rightline=false,
	nobreak,
	singleextra={%
		\draw(P-|O)++(-0.5em,0)node(tmp1){};
		\draw(P-|O)++(0.5em,0)node(tmp2){};
		\fill[black,rotate around={45:(P-|O)}](tmp1)rectangle(tmp2);
		\node at(P-|O){\color{white}\scriptsize\bf !};
		\draw[very thick](P-|O)++(0,-1em)--(O);%--(O-|P);
	}
}

% Define a custom environment for warning text
\newenvironment{warn}[1][Warning:]{ % Set the default warning to "Warning:"
	\medskip
	\begin{mdframed}[style=warning]
		\noindent{\textbf{#1}}
}{
	\end{mdframed}
}

%----------------------------------------------------------------------------------------
%	INFORMATION ENVIRONMENT
%----------------------------------------------------------------------------------------

% Usage:
% \begin{info}[optional title, defaults to "Info:"]
% 	contents
% 	\end{info}

\mdfdefinestyle{info}{%
	topline=false, bottomline=false,
	leftline=false, rightline=false,
	nobreak,
	singleextra={%
		\fill[black](P-|O)circle[radius=0.4em];
		\node at(P-|O){\color{white}\scriptsize\bf i};
		\draw[very thick](P-|O)++(0,-0.8em)--(O);%--(O-|P);
	}
}

% Define a custom environment for information
\newenvironment{info}[1][Info:]{ % Set the default title to "Info:"
	\medskip
	\begin{mdframed}[style=info]
		\noindent{\textbf{#1}}
}{
	\end{mdframed}
}
 % Include the file specifying the document structure and custom commands

%----------------------------------------------------------------------------------------
%	ASSIGNMENT INFORMATION
%----------------------------------------------------------------------------------------


\title{Análise de um Experimento Bem Planejado \\ \large Tarefa 7 --- Artigo Relacionado à Pesquisa} % Title of the assignment

\author{Giovani Tavares\\ \texttt{giovanitavares@usp.br}} % Author name and email address

\date{Universidade de São Paulo --- \today} % University, school and/or department name(s) and a date

%----------------------------------------------------------------------------------------

\begin{document}
	
	\maketitle % Print the title
	
	%----------------------------------------------------------------------------------------
	%	NOTA
	%----------------------------------------------------------------------------------------
		
	% =====================================================================
	\section{Artigo indicado}
	% =====================================================================
	
	\begin{itemize}
		\item \textbf{Título:} OG-RAG: Ontology-Grounded Retrieval-Augmented Generation for Large Language Models.
		\item \textbf{Autores:} Kartik Sharma, Peeyush Kumar e Yunqing Li.
		\item \textbf{Publicação:} Proceedings of the 2025 Conference on Empirical Methods in Natural Language Processing (EMNLP 2025), Suzhou, China, páginas 32962--32981. Association for Computational Linguistics.
		\item \textbf{DOI:} \texttt{10.18653/v1/2025.emnlp-main.1674}.
	\end{itemize}
	
	\noindent Este é um artigo fortemente relacionado à minha pesquisa de mestrado,
	que investiga o efeito da inclusão de uma camada de representação simbólica do
	conhecimento (ontologias) em arquiteturas de RAG, aplicada ao domínio de modelos
	metabólicos em escala genômica. O OG-RAG é justamente a arquitetura central que
	pretendo instanciar e avaliar, motivo pelo qual a análise de seu experimento é
	diretamente informativa para o desenho do meu próprio estudo.
	
	% =====================================================================
	\section{Descrição do experimento}
	% =====================================================================
	
	O experimento avalia se ancorar a recuperação de informação em uma ontologia de
	domínio melhora a qualidade das respostas geradas por LLMs, em comparação com
	métodos de recuperação sem estrutura simbólica. A arquitetura proposta representa
	os documentos de domínio como um \emph{hipergrafo}, no qual cada hiperaresta
	agrupa fatos relacionados segundo a ontologia, e recupera um conjunto mínimo de
	hiperarestas relevantes para cada consulta por meio de um algoritmo de otimização.
	
	A avaliação é conduzida sobre domínios especializados (como agricultura e
	trabalho baseado em conhecimento) e compara o OG-RAG a \emph{baselines} de
	recuperação. Os autores reportam que o OG-RAG aumenta em 55\% o \textit{recall}
	de fatos corretos e em 40\% a correção das respostas ao longo de quatro LLMs
	diferentes; adicionalmente, possibilita atribuição das respostas ao contexto 30\%
	mais rápida e melhora em 27\% a acurácia de raciocínio baseado em fatos em relação
	aos métodos de referência.
	
	% =====================================================================
	\section{Por que considero o experimento bem planejado e executado}
	% =====================================================================
	
	\begin{enumerate}
		\item \textbf{Isolamento da variável de interesse com \textit{baselines} claros.}
		O experimento compara a arquitetura proposta contra métodos de recuperação de
		referência, permitindo atribuir os ganhos observados à camada ontológica, e não
		a outros fatores do \textit{pipeline}. Esse cuidado com o \textit{baseline} é
		precisamente o tipo de controle que pretendo replicar na minha pesquisa.
		
		\item \textbf{Generalização entre modelos geradores.} As métricas são reportadas
		ao longo de \emph{quatro LLMs diferentes}, o que reduz o risco de que os ganhos
		sejam um artefato de um modelo específico e fortalece a validade externa das
		conclusões dentro dos domínios avaliados.
		
		\item \textbf{Métricas múltiplas e alinhadas ao objetivo.} A avaliação não se
		apoia em uma única métrica, mas em um conjunto que cobre \textit{recall} de
		fatos, correção das respostas, velocidade de atribuição ao contexto e raciocínio
		factual. Essa multiplicidade caracteriza melhor o comportamento do sistema do que
		um número isolado e é coerente com o objetivo declarado de recuperação
		fact-cêntrica.
		
		\item \textbf{Reprodutibilidade.} Os autores disponibilizam publicamente o
		código do método, o que viabiliza a verificação independente dos resultados e o
		reaproveitamento da implementação --- aspecto essencial para que eu possa adaptar
		a arquitetura ao meu domínio de aplicação.
		
		\item \textbf{Rastreabilidade como resultado mensurável.} O experimento trata a
		atribuição das respostas ao contexto como uma propriedade mensurável (e não apenas
		qualitativa), reforçando o argumento de que a estrutura simbólica melhora a
		auditabilidade das respostas, um ponto central da minha investigação.
	\end{enumerate}
	
	Em conjunto, esses elementos --- controle experimental por meio de
	\textit{baselines}, avaliação sobre múltiplos modelos, uso de várias métricas
	pertinentes e compromisso com a reprodutibilidade --- justificam minha avaliação
	de que se trata de um experimento bem planejado e executado, além de servir como
	referência metodológica direta para o desenho do meu próprio estudo comparativo.
	
	% =====================================================================
	\begin{thebibliography}{1}
		
		\bibitem{sharma2025ograg}
		Kartik Sharma, Peeyush Kumar, and Yunqing Li.
		\newblock OG-RAG: Ontology-grounded retrieval-augmented generation for large language models.
		\newblock In \textit{Proceedings of the 2025 Conference on Empirical Methods in Natural Language Processing (EMNLP 2025)}, pages 32962--32981, Suzhou, China, 2025. Association for Computational Linguistics.
		\newblock DOI: \texttt{10.18653/v1/2025.emnlp-main.1674}.
		
	\end{thebibliography}
	
\end{document}